\SourceLeaf{254}%
\Needspace{7\baselineskip}
\begin{center}
\textsc{Gauss} to \textsc{Schumacher}. Göttingen, 8 February 1846.
\end{center}

\mbox{.\ .\ .\ .\ .\ .\ .} The difference between right and left, however, can only be shown and cannot be defined, so it is similar in this respect to sweet and bitter. Omne simile claudicat, however. The latter applies only to beings that have organs of taste, the former to all minds capable of apprehending the material world. Yet two such minds can communicate directly about right and left only if one and the same individual material object forms a bridge between them. I say directly because A can also communicate with Z by having a material bridge established between A and B, another between B and C, and so on, or such bridges may already have been established. I stated succinctly in the Göttingischen Gelehrten Anzeigen 1831, p. 635, what bearing this matter has on metaphysics and that I find in it the decisive refutation of \textsc{Kant}'s fancy that space is MERELY the form of our external intuition. \mbox{.\ .\ .\ .\ .\ .\ .}

\par\bigskip

\Needspace{7\baselineskip}
\begin{center}
[\textsc{Gerling} \emph{to} \textsc{Gauss}. \emph{Marburg, 20 June 1846.}]
\end{center}

\{\mbox{.\ .\ .\ .\ .\ .\ .\ .} For the first time I found myself having to define a \textls{right}-handed screw for print [in the new edition of “Lorenz”]. I can think of no better way to put it than this: it is wound so that every point of the vertically standing screw off the axis can descend only if it goes from east through south to west. --- If there is a better definition, you would oblige me greatly by communicating it. \mbox{.\ .\ .\ .\ .\ .\ .}\}

\par\bigskip
